\documentclass{article}
\usepackage{graphicx} % Required for inserting images
\usepackage{datetime} % Add this line to include the datetime package
\usepackage{amsmath}
\usepackage{amssymb}
\usepackage{amsthm}
\usepackage{tikz}
%\usepackage{hyperref}
\usepackage{xcolor} % Required for defining custom colors

\definecolor{darkblue}{rgb}{0.0, 0.0, 0.5} % Define a darker blue color

\usepackage[colorlinks=true, linkcolor=darkblue, urlcolor=darkblue, citecolor=darkblue]{hyperref} % Use the darker blue color

% Define theorem style
\theoremstyle{plain} % Bold title, italic body
\newtheorem{theorem}{Theorem}
\newtheorem{lemma}[theorem]{Lemma}
\newtheorem{corollary}[theorem]{Corollary}

\theoremstyle{definition} % Bold title, normal body
\newtheorem{definition}[theorem]{Definition}
\newtheorem{example}[theorem]{Example}


\title{Homework 7}
\author{Aaron Honculada}
\date{\monthname[\month] \the\year}

\begin{document}

\maketitle

\section*{Problem 1}
\noindent \textit{Consider the Dihedral Group $D_n = \{r_0, r_1,... , r_{n-1}, s_0, ... , s_{n-1} \}$
with 
\begin{center}
    $r_ar_b = r_{a+b}, \quad r_as_b = s_{a+b}, \quad s_ar_b = s_{a-b}, \quad s_as_b = r_{a-b},$
\end{center}
where indices are taken modulo $n$. Describe the conjugacy classes of $D_n$ for all $n \geq 3$.}

We describe the conjugacy classes of $D_n$ in two cases: $n$ is even and $n$ is odd. 

When $n$ is even, the conjugacy classes can be described as follows:
\begin{itemize}
    \item $\{e\}$
    \item $\{r^k, r^{-k}\}$ for all $1 \leq k < \frac{n}{2}$
    \item $\{r^{\frac{n}{2}}\}$
    \item $\{r^{k}s\}$ for all even $0 \leq k < n$
    \item $\{r^{k}s\}$ for all odd $1 \leq k < n$.
\end{itemize}
Which gives a total of $1 + \frac{n}{2} - 1 + 1 + 2 = \frac{n}{2} + 3$ conjugacy classes.

In the case where $n$ is odd, the conjugacy classes are described as follows:
\begin{itemize}
    \item $\{e\}$
    \item $\{r^k, r^{-k}\}$ for all $1 \leq k \leq \frac{n-1}{2}$
    \item $\{r^ks\}$ for $0 \leq k \leq n-1$
\end{itemize}
Which gives a total of $1 + (\frac{n-1}{2}) + 1 = \frac{n+3}{2}$ conjugacy classes. 

\section*{Problem 2}
\noindent \textit{Let $G$ be a group acting on a set $X$. Show that if $x,y \in X$ lie in the same orbit,
then $Stab_x$ and $Stab_y$ are conjugate subgroups.}

\begin{proof}
    Given $x,y \in X$, since $x$ and $y$ lie in the same orbit, we have that
    there exists some $g \in G$ such that $x\cdot g = y$. $\text{Stab}_x$ and
    $\text{Stab}_y$ are conjugate subgroups if $\text{Stab}_y = g \cdot \text{Stab}_x \cdot g^{-1}$.  
    For $f \in \text{Stab}_y$ we have $y = f \cdot y \implies (g \cdot x) = f \cdot (g \cdot x) \implies (g \cdot x) =  (f \cdot g) \cdot x
    \implies g^{-1}g \cdot x = g^{-1}fg \cdot x \implies x = g^{-1}fg \cdot \implies g^{-1}fg \in \text{Stab}_x$ and
    $f \in g\text{Stab}_xg^{-1}$.
    Also for $f \in g\text{Stab}_xg^{-1}$ and $s \in  \text{Stab}_x$ we have $f = gsg^{-1}$.
    It follows that $f \cdot y = gsg^{-1} \cdot y = gsg^{-1} \cdot (gx) = gs \cdot x = g \cdot x = y \implies f \cdot y = y$. Therefore $f \in \text{Stab}_y$.

    We have shown that if $x,y \in X$ lie in the same orbit then $g\text{Stab}_xg^{-1} = \text{Stab}_y$ and $\text{Stab}_x$ and
    $\text{Stab}_y$ are conjugate subgroups.
\end{proof}

\section*{Problem 3}
\noindent \textit{Let $G$ be a finite simple group, and let $H$ be a proper subgroup
of $G$ of index $k$. Show that $|G|$ divides $k!$.}

\begin{proof}
    Given that $H$ is a proper subgroup of finite simple group $G$, and that $[G:H]=k$, we know that there
    are $k$ number of distinct left cosets of $H$ in $G$. Let $C$ be the set of these left cosets, $g'H$.
    We consider the action of $G$ on this set of cosets $C$. By definition
    that means there exists a group homomorphism $\varphi: G \to S_k$ since
    the elements of $g \in G$ permute the elements of $C$. Examining the kernel of $\varphi$ we
    have that ker$(\varphi)=\{g \in G \mid gg'H = g'H \text{ for all } g'H \in C \}$. By
    defining the kernel of $\varphi$ as such, we also have it that if $gg'H = g'H$ then
    $g'H = gg'H \implies (g')^{-1}gg' \in H$ and $g \in g'H(g')^{-1}$. So for all
    $g' \in G$, $g$ is in $\bigcap\limits_{g' \in G} g'Hg'^{-1}$. This is the definition of
    a normal subgroup of G. However we are given that $G$ is simple, therefore
    ker$(\varphi) = \{ e \}$ or $G$. If ker$(\varphi) = G$ then $gH = H$ which
    implies $H=G$. This contradicts our premise that $H$ is a proper subgroup of $G$.
    Therefore we must have it that ker$(\varphi) = \{ e\}$. Given that the kernel of $\varphi$
    is the trivial group, we know that $\varphi$ is injective and that $G \hookrightarrow S_k$.
    Therefore $|G| \mid |S_k|$ and $|G| \mid k!$.
\end{proof}

\section*{Problem 4}
\noindent \textit{Let $C$ be a conjugacy class of $S_n$. Show that one of the following
holds:
\begin{itemize}
    \item $C$ is disjoint from $A_n$
    \item $C$ is a conjugacy class of $A_n$
    \item $C$ splits into two conjugacy classes of $A_n$ of equal size
\end{itemize}
}

\section*{Problem 5}
\noindent \textit{Let $G$ be a finite group of order $n$, and let $p$ be a prime
dividing $n$. Let $X$ be the set of $p$-tuples $(g_1, g_2, ... , g_p)$, where
$g_i \in G$ for all $i$, and $g_1g_2 \cdots g_p = e$.}

\noindent (a) What is $|X|$?

For any $p$-tuple in $X$, when considering the first $p-1$ spots, at each spot there are $n$
possible elements to choose from. The final $p^{th}$ spot in the tuple must be the inverse of $g_1g_2...g_{p-1}$.
Therefore $|X| = n^{p-1}$. 

\noindent (b) Show that $\mathbb{Z}_p$ acts on $X$ by cyclic shift,
and describe the fixed points $X^{\mathbb{Z}_p}$ under this action.

$\mathbb{Z}_p$ acts on $X$ since for a given cyclic shift $g,h \in \mathbb{Z}_p$ and
for a $p$-tuple $x \in X$, if $g=0$, then $g\cdot x= x$ and for non-zero $g,h$ we have
$(g+h) \cdot x = g \cdot (h \cdot x)$. This is due to the observation that
cyclic shifts do not change the overall product of the elements of a given tuple.

The fixed points of $X^{\mathbb{Z}_p}$ only occur when
$g \cdot x = x$ which is when cycling the order of the elements of the tuple
does not change the product. This is only true when all of the elements of the tuple
are the same. Each fixed point in $X^{\mathbb{Z}_p}$ is of the form 
$(g,g,...,g)$ where for any pair in the tuple, $g_i = g_j$ for all $1 \leq i < j \leq p$ and
$g^p = e$.

\noindent (c) From parts (a) and (b), deduce that $G$ contains an element
of order $p$.

\begin{proof}
From the orbit-stabilizer theorem, we know that the size of an orbit of $X$ must
be equal to $\frac{|\mathbb{Z}_p|}{|Stab_x|} = \frac{p}{|Stab_x|}$.
Therefore every orbit under $\mathbb{Z}_p$ has either size 1 if it is a fixed point and the
stabilizer is all of $\mathbb{Z}_p$, or size $p$ if only
the identity is in the stabilizer. We know that there
is at least one fixed point since every element in the $p$-tuple can be the identity. We let $f =$ the number of fixed points
of $X$ under the action of $G$ and $f \geq 1$. So far we have $|X| = n^{p-1} \implies |X| - f = n^{p-1} - f$.
We also know the remaining elements of $X$ each lie in an orbit of size $p$. So for some $r$ amount
of orbits we have $|X| = n^{p-1} = f + rp$, and since we are
given $p \mid n$ this further imples that $f \equiv 0 \pmod{p} \implies p \mid f$
and $f \geq p$ since $f \geq 1$. We have shown that $G$ contains at least one
element of order $p$.
\end{proof}

\end{document}